\section{Experimental Results}
\label{sec:results}

\subsection{End-to-End AgentDojo Performance}

We first report all 24 fixed experimental conditions per model and then the stress-conditioned subset defined in Section~\ref{sec:stress-gate}. A condition is one attack framing, injection task, and environment suite. Results are macro averages: each condition receives equal weight even though banking, Slack, and workspace conditions contain six, two, and one user-task episodes, respectively.

\begin{figure*}[t]
\centering
\includegraphics[width=0.92\textwidth]{fig_end_to_end_frontier.pdf}
\caption{Macro-averaged utility and security on the stress-conditioned AgentDojo sets. Each point is the equal-weight mean over 18 DeepSeek or 16 Llama experimental conditions. Condition-bootstrap intervals are reported in Table~\ref{tab:main-results-revised}.}
\label{fig:end-to-end-frontier}
\end{figure*}

\begin{table*}[t]
\centering
\small
\setlength{\tabcolsep}{5pt}
\caption{Stress-conditioned AgentDojo results. Utility and security are macro means over conditions; brackets show 95\% condition-bootstrap intervals. Relative utility is the macro utility divided by the no-defense macro utility.}
\label{tab:main-results-revised}
\begin{tabular}{llrrr}
\toprule
Model & Defense & Utility [95\% CI] & Security [95\% CI] & Rel.\ utility \\
\midrule
DeepSeek-Chat (V3)
 & No defense & 67.6 [49.1, 84.3] & 38.0 [23.1, 52.8] & 1.00 \\
 & Spotlighting & 73.1 [56.5, 88.0] & 53.7 [38.0, 69.4] & 1.08 \\
 & Learned detector & 31.5 [23.1, 39.8] & 75.9 [54.6, 93.5] & 0.47 \\
 & \textbf{BioHarness} & \textbf{92.6 [88.9, 96.3]} & \textbf{98.1 [95.4, 100.0]} & \textbf{1.37} \\
\midrule
Llama-3-70B
 & No defense & 47.9 [31.2, 64.6] & 39.6 [27.1, 51.0] & 1.00 \\
 & Spotlighting & 52.1 [35.4, 69.8] & 53.1 [37.5, 67.7] & 1.09 \\
 & Learned detector & 52.1 [35.4, 69.8] & 70.8 [54.2, 85.4] & 1.09 \\
 & \textbf{BioHarness} & \textbf{51.0 [36.5, 66.7]} & \textbf{95.8 [89.6, 100.0]} & \textbf{1.07} \\
\bottomrule
\end{tabular}
\end{table*}

On DeepSeek, BioHarness raises macro security from 38.0\% without defense and 53.7\% with Spotlighting to 98.1\%, while raising macro utility from 67.6\% to 92.6\%. Both security and utility intervals separate cleanly from the no-defense and Spotlighting estimates. The learned detector reaches 75.9\% security but only 31.5\% utility, confirming that aggressive input blocking destroys task completion.

On Llama, BioHarness reaches 95.8\% macro security with 51.0\% utility---matching Spotlighting and the learned detector on utility while dominating them on security. Retention is 1.07 relative to the undefended baseline. Absolute utility tracks the open-weight model's task headroom: remasured banking utility stays near the no-defense banking rate, unlike DeepSeek where elastic cleavage recovered banking completion.

\paragraph{Unfiltered conditions.}
Table~\ref{tab:unfiltered-results} includes all 24 conditions. BioHarness remains at 98.6\% security on DeepSeek and 97.2\% on Llama. Including naturally failed attacks raises the no-defense security baseline, which is why stress conditioning matters for fair comparison.

\begin{table}[t]
\centering
\small
\setlength{\tabcolsep}{4pt}
\caption{Macro results over all 24 fixed conditions per model, without stress conditioning.}
\label{tab:unfiltered-results}
\begin{tabular}{llrr}
\toprule
Model & Defense & Utility & Security \\
\midrule
DeepSeek & No defense & 75.7 & 53.5 \\
 & Spotlighting & 75.7 & 61.1 \\
 & Learned detector & 40.3 & 77.8 \\
 & \textbf{BioHarness} & \textbf{94.4} & \textbf{98.6} \\
\midrule
Llama & No defense & 46.5 & 59.7 \\
 & Spotlighting & 51.4 & 58.3 \\
 & Learned detector & 59.7 & 80.6 \\
 & \textbf{BioHarness} & \textbf{54.9} & \textbf{97.2} \\
\bottomrule
\end{tabular}
\end{table}

\paragraph{Sensitivity to aggregation.}
Table~\ref{tab:weighted-results} weights every user-task episode equally. Security remains high under this alternative. On stress-conditioned DeepSeek episodes, BioHarness completes 87.5\% utility compared with 79.7\% without defense, so the utility gain survives the alternative estimand even though larger banking conditions receive more weight. Security stays near ceiling (96.9\%).

\begin{table}[t]
\centering
\small
\setlength{\tabcolsep}{4pt}
\caption{Episode-weighted sensitivity analysis on the stress-conditioned sets. DeepSeek contains 64 and Llama 61 user-task episodes per defense.}
\label{tab:weighted-results}
\begin{tabular}{llrr}
\toprule
Model & Defense & Utility & Security \\
\midrule
DeepSeek & No defense & 79.7 & 51.6 \\
 & Spotlighting & 81.2 & 64.1 \\
 & Learned detector & 34.4 & 90.6 \\
 & \textbf{BioHarness} & \textbf{87.5} & \textbf{96.9} \\
\midrule
Llama & No defense & 45.9 & 49.2 \\
 & Spotlighting & 44.3 & 59.0 \\
 & Learned detector & 44.3 & 77.0 \\
 & \textbf{BioHarness} & \textbf{45.9} & \textbf{93.4} \\
\bottomrule
\end{tabular}
\end{table}

\subsection{Results by Attack Framing}

Figure~\ref{fig:attack-type-tradeoff} separates four attack framings: direct injections in external content, explicit instruction overrides, education-framed attacks, and persona-framed attacks. The DeepSeek subgroup contains $(2,5,5,6)$ conditions and the Llama subgroup $(3,4,4,5)$ conditions.

\begin{figure*}[t]
\centering
\includegraphics[width=0.92\textwidth]{fig_attack_type_tradeoff.pdf}
\caption{Macro utility and security of BioHarness by attack framing on stress-conditioned conditions. Each point is the mean within one category. Identical points are offset slightly for visibility. Direct-injection estimates rest on only two DeepSeek and three Llama conditions.}
\label{fig:attack-type-tradeoff}
\end{figure*}

On DeepSeek, instruction-override, education-framed, and persona-framed conditions reach 100\% macro security with 93.3--94.4\% utility. The two direct-injection conditions average 83.3\% security and 83.3\% utility: residual misses concentrate where the adversary's instructions appear inside routine tool output without an explicit override frame, so the agent must rely on boundary checks rather than obvious goal conflict. On Llama, instruction/education/persona families remain at 100\% security, while the three direct-injection conditions average 77.8\% security; utility ranges from 43.3\% (persona) to 70.8\% (education).

\subsection{Controlled Repair Analysis}

The interactive sweeps measure full task completion. To isolate \emph{why} the harness preserves utility after an injection, we fix one injected workspace observation and evaluate five task-preservation checks together with one unauthorized email action. This probe measures mechanism-level recovery, not population AgentDojo success.

Figure~\ref{fig:utility-recovery} shows the progression. Withholding the complete suspicious output preserves security but satisfies only 20\% of the five utility checks. Removing the injection envelope while retaining benign metadata raises this score to 60\%. Allowing benign read-like continuation raises it to 80\%. The full goal-aligned authorization rule reaches 100\%. The unauthorized email remains blocked at every stage.

\begin{figure*}[t]
\centering
\includegraphics[width=0.92\textwidth]{fig_staged_utility_recovery.pdf}
\caption{Utility checks satisfied in one controlled injected-output replay. Scores change in 20-point increments because the case contains five utility checks. Security denotes one blocked malicious action and remains 100\% throughout.}
\label{fig:utility-recovery}
\end{figure*}

Removing metadata repair from the complete protocol reduces the case-study utility score by 40 points, and removing goal-aligned continuation reduces it by 20 points (Table~\ref{tab:module-ablation}). Five other single-component removals leave this particular case unchanged. A four-fixture extension covering file sharing, email, calendar creation, and hotel reservation satisfies all encoded utility checks and blocks the designated malicious action in all four fixtures.

\begin{table*}[t]
\centering
\small
\caption{Single-component removals on the controlled workspace replay. Utility differences are relative to the complete protocol; security is one blocked malicious action.}
\label{tab:module-ablation}
\begin{tabular}{lrr}
\toprule
Removed component & $\Delta U$ (points) & Security \\
\midrule
Watch-level read allowance & 0 & 100 \\
Lightweight injection state & 0 & 100 \\
Unified risk combination & 0 & 100 \\
Per-tool filtering & 0 & 100 \\
Goal-aligned continuation & $-20$ & 100 \\
Task-fidelity hint & 0 & 100 \\
Metadata repair & $-40$ & 100 \\
\bottomrule
\end{tabular}
\end{table*}

\subsection{Live Component Ablation}

The controlled probe suggests that metadata repair and goal-aligned continuation carry most of the utility recovery. We next ask which modules are necessary under the live AgentDojo protocol used in the main tables. Table~\ref{tab:live-dropone} reports macro means over all 24 DeepSeek conditions after removing one component from full BioHarness.

\begin{table}[t]
\centering
\small
\setlength{\tabcolsep}{4pt}
\caption{Live AgentDojo drop-one ablation on DeepSeek (macro mean over 24 conditions). Each row disables one module of BioHarness.}
\label{tab:live-dropone}
\begin{tabular}{lrr}
\toprule
Configuration & Utility & Security \\
\midrule
BioHarness (full) & 94.4 & 95.8 \\
$-$ reflection probe & 79.2 & 97.9 \\
$-$ feature-attribution reorder & 94.4 & 96.5 \\
$-$ path-consistency rules & 94.4 & 94.4 \\
$-$ goal-aligned continuity & 93.8 & 98.6 \\
$-$ metadata repair & 56.9 & 98.6 \\
\bottomrule
\end{tabular}
\end{table}

Metadata repair remains the dominant utility module: removing it drops macro utility by 37.5 points (94.4\%$\to$56.9\%) while security stays near ceiling. Workspace conditions fall from 100\% to 25\% utility without repair, matching the cargo-preserving cleavage design. Feature-attribution reordering and goal-aligned continuity leave utility essentially unchanged; path-consistency rules show a modest security cost (94.4\% vs.\ 95.8\% full).

The reflection probe is best read as a soft-control signal rather than an ASR booster. Removing it lowers macro utility from 94.4\% to 79.2\% (banking 83.3\%$\to$50.0\%; Slack 100\%$\to$87.5\%) while security remains high. On the scripted suite, overall attack block rate stays 75\% with or without reflection, but the blocking layer shifts: with reflection enabled, 20 of 36 blocked attacks are attributed to the reflection path, whereas without it those decisions move onto feature-attribution reordering. This pattern matches the design intent---surface drift early and allow goal-aligned or read-like continuations to soften lockdown---rather than adding a hard rule that raises stop rate. A standalone probe that scores proposed actions without the full Gate--Cleave loop still separates benign from malicious proposals only weakly, so reflection should be evaluated as part of the continuity protocol, not as an independent classifier.

A paired scripted audit under the same BioHarness profile reinforces the detection-side view. With path-consistency rules enabled, BioHarness blocks 71.4\% of in-distribution and 80.0\% of held-out scripted attacks at 0\% benign false-positive rate; disabling those rules lowers these rates to 57.1\% and 60.0\%. The residual misses concentrate on structured field smuggling, path smuggling, and retry paraphrases.

\subsection{Scripted Reordering Check}

As an implementation check, text-level reordering rules block all six designed order-dependent constructions once added on top of structural validation (Figure~\ref{fig:shuffle-gain}). This verifies intended coverage of the text path; live BioHarness uses feature-attribution reordering instead.

\begin{figure*}[t]
\centering
\includegraphics[width=0.92\textwidth]{fig_shuffle_security_gain.pdf}
\caption{Additional blocked-sample coverage from adding text-reordering rules to structural validation ($n=20$ samples per category).}
\label{fig:shuffle-gain}
\end{figure*}
